\section{Thinking 与 Instruct 模式的合并难题}

\subsection{Qwen3 的混合模式实验}

2025 年初，Qwen 团队心中有一个宏大的愿景：理想系统应该统一 thinking 和 instruct 模式，支持可调节的推理努力（类似低/中/高推理设置），甚至能根据提示和上下文自动推断合适的推理量。Qwen3 是这一方向上最清晰的公开尝试——它引入了``混合思考模式''，支持 thinking 和 non-thinking 行为共存于一个模型家族，并描述了一个包含``thinking mode fusion''的四阶段后训练流水线。

\subsection{合并的核心困难：数据分布冲突}

\begin{warningbox}{不仅仅是模型兼容性问题}
当人们谈论合并 thinking 和 instruct 时，往往首先想到的是模型侧兼容性：一个 checkpoint 能否支持两种模式、一个 chat template 能否切换。但更深层的问题是\textbf{两种模式的数据分布和行为目标存在本质差异}。
\end{warningbox}

作者坦率地总结了两种行为画像之间的张力：

\begin{itemize}
    \item \textbf{强 Instruct 模型}：追求直接性、简洁性、格式合规、低延迟，适合高频企业任务（重写、标注、模板化支持、结构化提取等）
    \item \textbf{强 Thinking 模型}：在困难问题上花费更多 token、维持连贯的中间结构、探索替代路径、保留足够的内部计算以提升最终正确性
\end{itemize}

\textit{These two behavior profiles pull against each other.}（这两种行为画像相互拉扯。）如果合并数据未经仔细筛选，结果通常是两个方向都表现平庸。

\subsection{分离路线 vs 集成路线}

实际上，Qwen 团队在 Qwen3 之后的 2507 版本中选择了\textbf{分离路线}——发布了独立的 Instruct 和 Thinking 变体（30B 和 235B）。大量商业客户仍然需要高吞吐、低成本、高可控的 instruct 行为。

而 Anthropic 选择了相反方向。Claude 3.7 Sonnet 被定位为混合推理模型，用户可以选择普通回复或扩展思考，API 用户可以设置 thinking budget。Anthropic 明确表示他们认为推理应该是一种集成能力而非独立模型。GLM-4.5 和 DeepSeek V3.1 也走了类似的混合路线。

\begin{importantbox}{有机合并的关键标准}
真正成功的合并需要的是\textbf{推理努力的平滑光谱}（a smooth spectrum of reasoning effort）。模型应该能够表达多个级别的努力，并理想地在它们之间自适应选择。GPT 风格的 effort control 指向了这一方向：\textit{a policy over compute, rather than a binary switch.}（对计算的策略，而非一个二元开关。）
\end{importantbox}

\subsection{本章小结}

Thinking 与 Instruct 的合并在概念上是对的，但数据分布冲突使其在工程上极具挑战。行业分化为分离路线（Qwen 2507）和集成路线（Anthropic、DeepSeek）。关键判断标准是合并是否``有机''——即是否实现了真正的推理努力连续谱。


